% Displacement-weighted jet substructure observables — draft
% Target: SciPost Physics / PRD (~15 pages). Plain article class for the
% draft; port to the journal class at submission.
% Build: pixi run paper   (tectonic paper/main.tex)
\documentclass[11pt]{article}

\usepackage[margin=2.5cm]{geometry}
\usepackage{amsmath,amssymb}
\usepackage{graphicx}
\usepackage{xcolor}
\usepackage{hyperref}
\usepackage{subcaption}

\graphicspath{{../plots/}}

% ------------------------------------------------------------------ macros
\newcommand{\stub}[1]{\begin{center}\fbox{\parbox{0.92\textwidth}{\color{red!70!black}\textbf{STUB:} #1}}\end{center}}
\newcommand{\todo}[1]{{\color{red!70!black}[TODO: #1]}}
\newcommand{\dnought}{d_0}
\newcommand{\ctau}{c\tau}
\newcommand{\pid}{\pi_{\mathrm{D}}}
\newcommand{\qD}{q_{\mathrm{D}}}
\newcommand{\dR}{\Delta R}
\newcommand{\pT}{p_{\mathrm{T}}}

\title{Displacement-weighted jet substructure observables\\
for emerging-jet signatures}
\author{Jackson Burzynski\thanks{\texttt{jackson.carl.burzynski@cern.ch}}
\and \todo{student} \and \todo{optional theorist}}
\date{\today\ --- \textbf{DRAFT}}

\begin{document}
\maketitle

\begin{abstract}
Dark-sector models with confining hidden dynamics predict ``emerging
jets'': QCD-like jets whose constituents originate from displaced
dark-pion decays. Model-agnostic anomaly-detection searches built on
jet substructure --- autoencoders and related unsupervised
architectures trained on calorimeter- or momentum-based observables ---
are largely \emph{blind} to this signature: we show that the standard
substructure basis (angularities, energy correlation functions,
$N$-subjettiness) carries essentially no lifetime information, so an
anomalous sector whose only distinguishing feature is displacement
would evade substructure-based anomaly detection.
We introduce an interpretable, analytic family of
\emph{displacement-weighted} jet substructure observables ---
$\dnought$-weighted angularities, displaced energy--energy correlators,
and lifetime moments --- and characterize their discrimination power as
a function of the dark-pion proper lifetime,
$\ctau(\pid) \in [1, 300]$~mm, against light-flavor \emph{and
heavy-flavor} QCD jets. The pairwise correlators, which exploit the
angular correlation of displacement among tracks from a common decay
vertex, dominate the single-track observables precisely in the
short-lifetime regime where dark pions must be distinguished from $B$
hadrons. Augmenting a (variational) autoencoder's substructure input basis with
these observables extends the generality of anomaly detection to
lifetime-frontier signatures: at a fixed $1\%$ QCD anomaly rate, the
signal efficiency rises from the background-consistent $\lesssim 2\%$
(standard basis, at every lifetime) to $54\%$ at $\ctau = 1$~mm and
$\geq 90\%$ for $\ctau \geq 10$~mm, with the QCD background model
unaffected. We further
quantify the fraction of the fully supervised discrimination ceiling,
defined by a per-lifetime boosted-decision-tree classifier on the
complete observable basis, that individual observables capture, and
demonstrate robustness under a parametrized tracking response
including impact-parameter resolution, efficiency loss with production
radius, large-radius tracking, and pileup.
\end{abstract}

\tableofcontents

% =====================================================================
\section{Introduction}
\label{sec:intro}

Confining dark sectors coupled to the Standard Model through a heavy
mediator generically produce collider signatures in which a parton-level
dark quark showers and hadronizes into a multiplicity of dark hadrons,
a fraction of which decay back to SM particles with macroscopic proper
lifetimes~\cite{Schwaller:2015gea,Albouy:2022cin}. The resulting
``emerging jet'' contains an anomalously large fraction of tracks with
sizable impact parameter, distributed among several displaced vertices
along the jet axis.

Existing search strategies occupy two extremes. Cut-based analyses use
simple per-jet counting variables --- e.g.\ the prompt-track energy
fraction or the median unsigned impact parameter --- which are robust
and interpretable but discard the rich correlation structure of the
displaced final state~\cite{Schwaller:2015gea,ATLAS:EJ,CMS:EJ}.
Machine-learned constituent-level taggers, at the other extreme,
approach the information-theoretic ceiling but resist physical
interpretation, complicate calibration, and obscure which features of
the signal carry the discrimination power.

A third, increasingly prominent strategy is \emph{model-agnostic
anomaly detection}: unsupervised or weakly supervised architectures ---
autoencoders (AEs), variational autoencoders (VAEs), and density-based
methods --- trained on background-dominated data to flag any jet or
event that populates the tails of the learned background
model~\cite{Farina:2018fyg,Heimel:2018mkt,Kasieczka:2021xcg,ATLAS:dijetAD}.
The generality of such searches is bounded by their input
representation: an anomaly can only be flagged if it is anomalous
\emph{in the variables the network sees}. Substructure-based anomaly
detection typically uses calorimeter- or momentum-based inputs ---
constituent four-vectors, angularities, energy correlators,
$N$-subjettiness --- none of which carry lifetime information. This
paper demonstrates that blind spot quantitatively: across every
standard substructure observable we study, emerging jets are separated
from QCD by rejection factors of only $2$--$5$, independent of the
dark-pion lifetime (Sec.~\ref{sec:results}). A dark shower whose
principal anomaly is \emph{displacement} looks ordinary to a
substructure autoencoder. The throughline of this work is that
augmenting anomaly-detection input bases with a small set of
interpretable, displacement-aware observables closes this blind spot
at negligible architectural cost (Sec.~\ref{sec:anomaly}).

This paper develops the middle ground: a systematic, analytic family of
jet substructure observables in which each track is weighted by its
displacement significance. The family is organized in three tiers of
increasing structure --- weighted angularities (single-track moments),
displaced energy--energy correlators (pairwise moments), and lifetime
moments (direct $\ctau$ estimators) --- and reduces smoothly to standard
track-based substructure in the prompt limit. Because track-level
observables are not infrared-and-collinear safe, a first-principles
description would proceed through track
functions~\cite{Chang:2013rca,Li:2021zcf}; we define our observables
with this correspondence in mind but defer calculability questions.
\todo{tighten IRC/track-function framing with theorist input}

Our physics goals are:
(i) define the observable family and identify which moments carry
lifetime information;
(ii) map discrimination power as a function of $\ctau(\pid)$ against
light-flavor and heavy-flavor QCD jets separately --- the irreducible
background at short lifetime being $b$ jets, since
$\ctau(B) \approx 0.5$~mm;
(iii) demonstrate that a (V)AE anomaly detector trained on a standard
substructure basis gains sensitivity to displaced signatures across the
full $\ctau$ range when the input basis is augmented with the
displacement-weighted observables, without degrading its background
model;
(iv) quantify how much of the fully supervised discrimination ceiling
on the complete observable basis individual observables retain; and
(v) establish robustness under a realistic parametrized tracking
response.

% =====================================================================
\section{Observable definitions}
\label{sec:observables}

Jets are clustered with the anti-$k_t$ algorithm~\cite{Cacciari:2008gp}
with $R = 1.0$ from all visible final-state particles. Observables are
computed from the jet's charged constituents (``tracks'') $i$ with
transverse-momentum fraction $z_i = p_{\mathrm{T},i} / \pT^{\mathrm{jet}}$, angle
$\theta_i = \dR(i, \mathrm{jet\ axis})/R$, transverse impact parameter
$d_{0,i}$, and per-track resolution $\sigma_i$.

\subsection{Displacement weight}
\label{sec:weight}

Each track carries a displacement weight based on its impact-parameter
\emph{significance},
\begin{equation}
  w_i \;=\; \ln\!\left(1 + \frac{|d_{0,i}|}{\sigma_i}\right).
  \label{eq:weight}
\end{equation}
The logarithm compresses the long lifetime tail, and the significance
ratio makes the weight self-regulating against resolution: a prompt
track smeared to $|\dnought| \sim \sigma$ receives $w \simeq \ln 2$,
while genuinely displaced tracks receive $w \sim \ln(|\dnought|/\sigma)
\gg 1$. In the ideal-detector limit $\sigma \to 0$ the weight diverges
only logarithmically, so all observables below remain finite and
smearing-stable.

Each family below is presented as a \emph{nominal} (lifetime-blind)
observable next to its \emph{displacement-weighted} counterpart,
obtained by a single insertion of the weight of Eq.~(\ref{eq:weight})
per track. This pairing is used throughout the paper: every
distribution, ROC, and anomaly-detection input basis is shown for both
members, so the lifetime information contributed by the weighting is
read off directly. Table~\ref{tab:observables} summarizes the family.

\subsection{Angularities}
\label{sec:tier1}

The nominal track angularities~\cite{Larkoski:2014pca} and their
displacement-weighted generalization are
\begin{equation}
  \lambda^\kappa_\beta \;=\; \sum_{i \in \mathrm{jet}}
      z_i^\kappa \, \theta_i^\beta ,
  \qquad\qquad
  \lambda^\kappa_\beta(w) \;=\; \sum_{i \in \mathrm{jet}}
      w_i \, z_i^\kappa \, \theta_i^\beta .
  \label{eq:angularity}
\end{equation}
The nominal family contains the standard substructure workhorses:
girth $(\kappa,\beta) = (1,1)$, the mass-like angularity $(1,2)$,
track multiplicity $n_{\mathrm{trk}} = \lambda^0_0$, and
$(\pT^D)^2 = \lambda^2_0$. The corresponding weighted anchors are the
displaced-$\pT$ fraction $(1,0)$, displaced girth $(1,1)$, and a
displaced-mass proxy $(1,2)$. In the fully prompt limit every track
sits at the resolution pedestal $w_i \simeq \ln 2$, so
$\lambda^\kappa_\beta(w) \to \ln 2 \cdot \lambda^\kappa_\beta$: the
weighted observable degenerates to a rescaled copy of its nominal
partner, and any \emph{shape} difference between the pair is a pure
lifetime effect. A scan in $(\kappa, \beta)$ identifies which moments
carry the lifetime information (Sec.~\ref{sec:kbscan}); it strongly
favors the corner $\lambda^0_0(w) = \sum_i w_i$, the
\emph{displacement-weighted track multiplicity}, which we promote to a
flagship observable.

\subsection{Energy--energy correlators}
\label{sec:tier2}

The nominal two-point energy correlation
function~\cite{Larkoski:2013eya} and its displaced counterpart are
\begin{equation}
  \mathrm{ECF}_2(\beta) \;=\; \sum_{i<j} z_i z_j \, (\dR_{ij})^\beta,
  \qquad\qquad
  \mathrm{dEEC}(\beta) \;=\; \sum_{i<j} z_i z_j \,
      (\dR_{ij})^\beta \; g(w_i, w_j),
  \label{eq:deec}
\end{equation}
with pair weighting $g = w_i w_j$ or $g = \min(w_i, w_j)$.
The physics carried by the weighting is qualitatively new at this
order: dark-pion decays produce \emph{pairs} of displaced tracks from
common vertices, so displacement in signal is angularly correlated at
small pair separation $\dR_{ij}$, while in QCD (including $B$ decays,
with their single low-multiplicity vertex chain) the correlation is
weaker. The product weighting rewards pairs in which both tracks are
displaced; the $\min$ weighting is stricter, vanishing unless
\emph{both} tracks are significantly displaced, and is correspondingly
more robust against single-track resolution tails. The differential
profile $\langle w_i w_j \rangle (\dR_{ij})$, shown in
Fig.~\ref{fig:wij}, exhibits the vertex pair correlation directly and
connects to the energy-correlator-on-tracks literature
\cite{Chang:2013rca,Lee:2023npz}.

The construction extends to higher point count. Writing
$\mathrm{dECF}_1 = \sum_i z_i w_i$ and
$\mathrm{dECF}_2 = \mathrm{dEEC}$, the displaced three-point correlator
is
\begin{equation}
  \mathrm{dECF}_3(\beta) \;=\; \sum_{i<j<k} z_i z_j z_k \,
      (\dR_{ij}\dR_{ik}\dR_{jk})^\beta \; g(w_i, w_j, w_k),
\end{equation}
alongside the nominal $\mathrm{ECF}_3$ and ratios
$C_2 = \mathrm{ECF}_3 \mathrm{ECF}_1 / \mathrm{ECF}_2^2$,
$D_2 = \mathrm{ECF}_3 \mathrm{ECF}_1^3 / \mathrm{ECF}_2^3$. With the
product pairing $g = w_i w_j w_k$, each $\mathrm{dECF}_n$ carries $n$
powers of the weight, so in the displaced ratios
\begin{equation}
  C_2(w) = \frac{\mathrm{dECF}_3\,\mathrm{dECF}_1}{\mathrm{dECF}_2^2},
  \qquad
  D_2(w) = \frac{\mathrm{dECF}_3\,\mathrm{dECF}_1^3}{\mathrm{dECF}_2^3},
\end{equation}
the powers of $w$ cancel exactly as the powers of $z$ do: for a jet in
which all tracks share a common weight, $C_2(w)$ and $D_2(w)$ reduce
identically to their nominal partners. They are therefore
\emph{displacement-normalized} shape variables --- insensitive to the
overall amount of displacement, and sensitive instead to its
angular and vertex \emph{structure}: a dark shower distributes
displacement over many vertices across the jet, while a $B$-decay
chain concentrates it in one.

\subsection{Lifetime moments}
\label{sec:tier3}

The lifetime moments have no nontrivial nominal partner --- the
$\dnought$-independent limit of the family is
$L_0 = \sum_i z_i$, the charged energy fraction --- and are the most
direct $\ctau$ estimators:
\begin{equation}
  L_n \;=\; \sum_i z_i \, |d_{0,i}|^n ,
  \qquad n = 0, 1, 2, \ldots
  \label{eq:moments}
\end{equation}
together with dimensionless ratios such as $L_2 L_0 / L_1^2$, which
measures the width of the per-jet displacement distribution
independently of its scale. These support a ``measure, don't just
tag'' program: given a tagged emerging-jet sample, the moments estimate
$\ctau(\pid)$ itself. \todo{develop the $\ctau$-measurement section}

\subsection{Signed variants}
\label{sec:signed}

Throughout, we also consider signed-$\dnought$ variants in which
$|d_{0,i}|$ in Eq.~(\ref{eq:weight}) carries the lifetime sign of the
track relative to the jet axis: positive if the point of closest
approach lies along the jet direction, negative otherwise. As in
heavy-flavor tagging, the negative-signed tail is populated purely by
resolution and calibrates the $w_i \simeq \ln 2$ prompt pedestal in
situ. \stub{Signed-variant distributions and in-situ calibration
closure test.}

\subsection{Prong-structure observables and search-style baselines}
\label{sec:standard}

Two further sets complete the basis. First, the nominal
prong-structure observables, for which we do not construct weighted
counterparts here: $N$-subjettiness~\cite{Thaler:2010tr} with
exclusive-$k_t$ axes,
\begin{equation}
  \tau_N = \frac{\sum_i z_i \, \min_a (\dR_{ia})^\beta}
                {\sum_i z_i \, R^\beta},
\end{equation}
its ratios $\tau_{21}$, $\tau_{32}$, and the jet mass. Second, the
search-style single-track baselines against which any new displaced
observable must be compared: the mean impact-parameter significance
$\langle |\dnought|/\sigma \rangle$ and the prompt $\pT$ fraction
(fraction of track $\pT$ carried by tracks with
$|\dnought|/\sigma < 3$), the latter analogous to the $\alpha$-type
variables of Refs.~\cite{Schwaller:2015gea,CMS:EJ}.

\begin{table}[tb]
  \centering
  \caption{The observable basis. Each row pairs a nominal
  (lifetime-blind) observable with its displacement-weighted
  counterpart; the union of the ``nominal'' column is the
  anomaly-detection input basis $S$, and adding the ``weighted'' column
  gives $S{+}D$ (Sec.~\ref{sec:anomaly}).}
  \label{tab:observables}
  \begin{tabular}{lll}
    \hline
    family & nominal & displacement-weighted \\
    \hline
    angularities
      & $\lambda^\kappa_\beta = \sum_i z_i^\kappa \theta_i^\beta$
      & $\lambda^\kappa_\beta(w) = \sum_i w_i z_i^\kappa \theta_i^\beta$ \\
    \quad anchors
      & girth, mass ang., $n_{\mathrm{trk}}$, $(\pT^D)^2$
      & displ.\ $\pT$ frac., displ.\ girth, displ.\ mass \\
    2-point correlator
      & $\mathrm{ECF}_2(\beta)$
      & $\mathrm{dEEC}(\beta)$, $g = w_i w_j$ or $\min(w_i, w_j)$ \\
    3-point correlators
      & $\mathrm{ECF}_3$, $C_2$, $D_2$
      & $\mathrm{dECF}_3$, $C_2(w)$, $D_2(w)$ \\
    prong structure
      & $\tau_1$, $\tau_{21}$, $\tau_{32}$, jet mass
      & --- \\
    lifetime moments
      & $L_0 = \sum_i z_i$ (trivial)
      & $L_n = \sum_i z_i |d_{0,i}|^n$, $L_2 L_0/L_1^2$ \\
    baselines
      & ---
      & $\langle |\dnought|/\sigma \rangle$, prompt $\pT$ fraction \\
    \hline
  \end{tabular}
\end{table}

% =====================================================================
\section{Simulation and event selection}
\label{sec:simulation}

\subsection{Signal and background samples}

Signal events are generated with \textsc{Pythia}~8.312
\cite{Bierlich:2022pfr} using the Hidden Valley
module~\cite{Carloni:2010tw}: $s$-channel production
$pp \to Z' \to \qD\bar{\qD}$ at $\sqrt{s} = 13.6$~TeV with
$m_{Z'} = 1.5$~TeV, followed by a dark shower and hadronization with
$N_{\mathrm{gauge}} = 3$, running dark coupling with
$\Lambda_{\mathrm{D}} = 10$~GeV, $m_{\qD} = 10$~GeV,
$m_{\pid} = 5$~GeV, $m_{\rho_{\mathrm{D}}} = 20$~GeV, and vector
fraction $0.75$, following the unflavored benchmark of the Dark Showers
Task Force~\cite{Albouy:2022cin} and the original emerging-jets
proposal~\cite{Schwaller:2015gea}. Dark rhos decay promptly to dark
pion pairs; dark pions decay to $d\bar d$ with proper lifetime scanned
over $\ctau(\pid) \in \{0, 1, 3, 10, 30, 100, 300\}$~mm, where $0$
denotes the prompt limit. One alternate mediator portal is deferred to
Appendix~\ref{app:portal}.

The background is QCD dijet production in the same $\hat p_T$ regime,
split by jet flavor at analysis level: a jet is labeled $b$ ($c$) if a
weakly-decaying bottom (charm) hadron with $\pT > 5$~GeV lies within
$\dR < 1.0$ of the jet axis. A dedicated hard-$b\bar b$ sample
augments the $b$-jet statistics; per-flavor results are shape-based, so
the samples are combined freely. \todo{hard-$c\bar c$ sample}

\subsection{Jet and track selection}

Jets: anti-$k_t$, $R=1.0$, built from all visible final-state
particles; the two leading jets with
$500 < \pT < 1000$~GeV and $|\eta| < 2.5$ are analyzed.
Tracks: charged constituents with $\pT > 1$~GeV, $|\eta| < 2.5$, and
truth $|\dnought| < 300$~mm. The truth impact parameter is computed
from the production vertex in the straight-line approximation,
$\dnought = (v_x p_y - v_y p_x)/\pT$. In all discrimination results the
background jets are reweighted, per flavor, so that their $\pT$
spectrum matches the (lifetime-independent) signal jet spectrum;
rejection factors therefore reflect shape differences in the
observables, not residual kinematic differences between the $Z'$ decay
and the QCD $\hat p_T$ slice.

\subsection{Parametrized tracking response}
\label{sec:tracking}

Detector effects are applied as a parametrized transform of the truth
track collection, so identical observable code runs on truth-level and
smeared inputs. The impact-parameter resolution is
\begin{equation}
  \sigma(\dnought) = a \oplus b/\pT, \qquad
  a = 10~\mu\mathrm{m}, \quad b = 70~\mu\mathrm{m}\,\mathrm{GeV},
\end{equation}
and $\sigma_i$ in Eq.~(\ref{eq:weight}) always refers to this
parametrization, including at truth level. The longitudinal impact
parameter $z_0$, computed at the transverse point of closest approach,
is smeared analogously with
$\sigma(z_0) = 20~\mu\mathrm{m} \oplus 100~\mu\mathrm{m}\,\mathrm{GeV}/\pT$. Reconstruction efficiency
falls with production radius $r$ as
$\varepsilon(r) = \varepsilon_0 / (1 + (r/r_{1/2})^2)$ with a hard
acceptance ceiling $r < r_{\max}$. Two scenarios bracket realistic
tracking: \emph{standard} ($r_{1/2} = 40$~mm, $r_{\max} = 120$~mm) and
\emph{standard + large-radius tracking}
($r_{1/2} = 150$~mm, $r_{\max} = 300$~mm).

\stub{Tracking-scenario results: observable degradation vs $\ctau$
under standard vs standard+LRT; this comparison is itself a physics
result. Samples and code exist; run \texttt{pixi run plots --scenario
standard / standard\_lrt}.}

Pileup is modeled by overlaying Poisson($\mu$) minimum-bias events
(\textsc{Pythia} \texttt{SoftQCD:inelastic}) per jet before the tracking
transform, with each overlaid vertex displaced longitudinally by the
beam-spot spread ($\sigma_z = 45$~mm); overlaid tracks receive the
identical smearing, efficiency, and selection as hard-scatter tracks.
Track selection follows LLP-search practice: a track is kept if it is
primary-vertex associated ($|z_0 \sin\theta| < 1.5$~mm) \emph{or} is a
displaced candidate ($|\dnought|/\sigma > 3$). Prompt pileup from other
vertices fails both conditions; the surviving contamination is the
genuinely displaced pileup component from strange-hadron decays.

% =====================================================================
\section{Truth-level results}
\label{sec:results}

\subsection{Distributions}

Figure~\ref{fig:dists} shows the two flagship observable pairs side by
side: nominal girth against displaced girth, and $\mathrm{ECF}_2$
against dEEC. In each nominal panel (left column) the signal
distributions at different lifetimes lie nearly on top of one another
and inside the QCD envelope; inserting the displacement weight (right
column) fans the same jets out by lifetime and separates every
$\ctau \geq 1$~mm point from all QCD flavors, with $b$ jets closest to
signal. This left-column/right-column contrast is the paper's central
construction: the columns differ by nothing but the per-track weight
of Eq.~(\ref{eq:weight}).

\begin{figure}[tb]
  \centering
  \includegraphics[width=0.49\textwidth]{dist_girth_truth.png}
  \includegraphics[width=0.49\textwidth]{dist_ang_11_truth.png}\\
  \includegraphics[width=0.49\textwidth]{dist_eec_b1_truth.png}
  \includegraphics[width=0.49\textwidth]{dist_deec_min_truth.png}
  \caption{Nominal observables (left column) and their
  displacement-weighted counterparts (right column), for QCD jets by
  flavor and signal at representative lifetimes, truth-level tracking.
  Top: girth $\lambda^1_1$ vs displaced girth $\lambda^1_1(w)$.
  Bottom: $\mathrm{ECF}_2(\beta{=}1)$ vs
  $\mathrm{dEEC}(\beta{=}1, \min)$.}
  \label{fig:dists}
\end{figure}

\begin{figure}[tb]
  \centering
  \includegraphics[width=0.62\textwidth]{wij_profile_truth.png}
  \caption{Energy-weighted profile of the pair displacement correlation
  $\langle w_i w_j\rangle$ versus pair separation $\dR_{ij}$. The
  small-$\dR$ enhancement in signal reflects track pairs from common
  dark-pion decay vertices; QCD $b$ jets show a much smaller
  enhancement from the single $B$-decay chain.}
  \label{fig:wij}
\end{figure}

\subsection{Discrimination power versus lifetime}

Figure~\ref{fig:rejection} shows the background rejection at 50\%
signal efficiency as a function of $\ctau(\pid)$, per background
flavor. Three qualitative results emerge:

\begin{enumerate}
  \item \textbf{Pairwise beats single-track at short lifetime.} At
  $\ctau = 3$~mm the dEEC observables reach $b$-jet rejections of
  $\mathcal{O}(300)$, versus $\mathcal{O}(5)$ for the single-track
  baselines $\langle|\dnought|/\sigma\rangle$ and $L_1$, and
  $\mathcal{O}(10\text{--}30)$ for the weighted angularities. Even at
  $\ctau = 1$~mm --- barely twice the $B$-hadron lifetime --- the dEEC
  retains $\mathcal{O}(100)$ rejection. The angular pair correlation,
  not the total displacement, is what distinguishes a 3~mm dark pion
  from a $B$ hadron. The three-point $\mathrm{dECF}_3$ sharpens the
  coincidence requirement further and becomes the strongest single
  observable for $\ctau \gtrsim 10$~mm (e.g.\ $b$-jet rejection
  $\sim 4\times 10^3$ at 30~mm, three times the dEEC), while the
  displacement-normalized
  ratio $C_2(w)$ shows a distinct and complementary behavior: its
  rejection is nearly \emph{independent of lifetime}
  ($\approx 20$--$25\times$ against $b$ jets for
  $1 \le \ctau \le 300$~mm), as expected for a variable sensitive to
  the vertex structure of the displacement rather than its amount ---
  a uniform working point across the entire lifetime grid.
  \item \textbf{Standard substructure carries no lifetime information.}
  All lifetime-blind observables sit at rejections of $2$--$5$,
  independent of $\ctau$, confirming that the discrimination above is
  attributable entirely to the displacement weighting. This is the
  quantitative statement of the anomaly-detection blind spot motivated
  in Sec.~\ref{sec:intro}: in any input basis built from these
  observables, emerging jets are at most mildly atypical, and no
  unsupervised method --- however expressive --- can recover
  information absent from its inputs.
  \item \textbf{The long-lifetime regime approaches saturation even at
  $10^5$-jet background samples.} With $8.6\times 10^4$ $b$ jets and
  $7.6\times 10^4$ light jets, most rejections are now measured rather
  than bounded --- e.g.\ $\Sigma_i w_i$ reaches a measured $b$-jet
  rejection of $2\times 10^5$ at $\ctau = 300$~mm --- while the
  prompt-$\pT$-fraction baseline still rejects every background jet
  for $\ctau \geq 10$~mm (statistics lower bounds, open markers).
  \todo{final production statistics on the cluster}
\end{enumerate}

\begin{figure}[tb]
  \centering
  \includegraphics[width=0.49\textwidth]{rejection_QCD_b_truth_eff50.png}
  \includegraphics[width=0.49\textwidth]{rejection_QCD_light_truth_eff50.png}
  \caption{Background rejection at 50\% signal efficiency versus
  $\ctau(\pid)$ for $b$ jets (left) and light jets (right).
  Solid: displacement-weighted observables; dashed: their nominal
  partners. Missing points indicate rejection exceeding the sample
  statistics. Truth-level tracking.}
  \label{fig:rejection}
\end{figure}

\subsection{Which moments carry the lifetime information?}
\label{sec:kbscan}

Figure~\ref{fig:kbscan} shows the $(\kappa,\beta)$ scan of the weighted
angularities: $b$-jet rejection at 50\% signal efficiency over the
grid $\kappa, \beta \in [0, 2]$. The result is a monotone gradient
toward the origin: the best moment is
$\lambda^0_0(w) = \sum_i w_i$, the displacement-weighted track
multiplicity, and every power of $z_i$ or $\theta_i$ \emph{dilutes} the
discrimination. Physically, the lifetime information resides in the
\emph{number} of significantly displaced tracks --- dark-pion decay
products are numerous but soft and broadly distributed within the jet,
so energy weighting suppresses exactly the tracks that carry the
signal. The single-observable hierarchy at short lifetime is therefore:
pairwise correlation (dEEC) $\gtrsim$ weighted multiplicity
$\sum_i w_i$ $\gg$ energy-weighted angularities $\gg$ single-track
means. This motivates the flagship set carried into
Secs.~\ref{sec:anomaly}--\ref{sec:gap}: $\{\mathrm{dEEC}(\min),\;
\sum_i w_i,\; \lambda^1_0(w)\}$.

\begin{figure}[tb]
  \centering
  \includegraphics[width=0.49\textwidth]{kb_scan_QCD_b_ctau3_truth.png}
  \includegraphics[width=0.49\textwidth]{kb_scan_QCD_b_ctau30_truth.png}
  \caption{$b$-jet rejection at 50\% signal efficiency for the
  displacement-weighted angularities $\lambda^\kappa_\beta(w)$ over the
  $(\kappa,\beta)$ grid, at $\ctau = 3$~mm (left) and 30~mm (right).
  ``$>$'' marks sample-statistics lower bounds.}
  \label{fig:kbscan}
\end{figure}

% =====================================================================
\section{Robustness under realistic tracking}
\label{sec:robustness}

Figure~\ref{fig:scenarios} compares the flagship observables across the
three tracking scenarios of Sec.~\ref{sec:tracking}. Three lessons
emerge. First, the pairwise dEEC pays an \emph{efficiency-squared}
penalty: under standard tracking its short-lifetime advantage largely
collapses (at $\ctau = 3$~mm the $b$-jet rejection falls from
$\mathcal{O}(300)$ at truth level to $\mathcal{O}(3)$), because a pair
observable requires both decay products to be reconstructed. Second,
large-radius tracking restores much of the pairwise information for
$\ctau \gtrsim 10$~mm, and the LRT-vs-standard gap as a function of
$\ctau$ directly quantifies the physics value of LRT for emerging-jet
searches. Third, counting-type observables (prompt $\pT$ fraction,
weighted multiplicity) are the most tracking-robust at short lifetime:
losing a displaced track degrades them linearly rather than
quadratically. A realistic search should therefore combine a
counting observable (robust) with the dEEC (maximal truth-level
information), and the ML studies of
Secs.~\ref{sec:anomaly}--\ref{sec:gap} are performed per scenario.
\todo{quantitative degradation table}

\begin{figure}[tb]
  \centering
  \includegraphics[width=0.62\textwidth]{scenario_comparison_QCD_b.png}
  \caption{$b$-jet rejection at 50\% signal efficiency versus
  $\ctau(\pid)$ for the flagship observables under the three tracking
  scenarios: truth (solid), standard tracking (dashed), and standard +
  large-radius tracking (dash-dotted). Open triangles are
  sample-statistics lower bounds.}
  \label{fig:scenarios}
\end{figure}

\paragraph{Pileup.} Figure~\ref{fig:pileup} shows the flagship
observables under the standard+LRT scenario with and without minimum-bias
overlay at $\mu = 60$. The rejection curves are statistically
indistinguishable across the full lifetime grid. The mechanism is
twofold: the $z_0$ association removes the prompt pileup bulk, and the
significance weighting caps what leaks through at the resolution
pedestal $w \simeq \ln 2$. After selection, pileup constitutes only
$\sim 4\%$ of signal-jet tracks at $\mu = 60$, and while roughly half
of that residual is genuinely displaced (strange-hadron daughters,
Fig.~\ref{fig:pileup} right), it is angularly uncorrelated and too rare
to populate displaced \emph{pairs}: the dEEC is entirely unaffected.
Displacement-weighted substructure is therefore pileup-robust by
construction, in contrast to track-counting variables whose
denominators grow with every soft track.

\begin{figure}[tb]
  \centering
  \includegraphics[width=0.49\textwidth]{pileup_rejection_QCD_b_standard_lrt.png}
  \includegraphics[width=0.49\textwidth]{pileup_composition_standard_lrt.png}
  \caption{Left: $b$-jet rejection versus $\ctau(\pid)$ for the flagship
  observables with (dotted, $\mu = 60$) and without (solid) minimum-bias
  pileup overlay, standard+LRT tracking. Right: post-selection
  $|\dnought|/\sigma$ of hard-scatter versus pileup tracks in signal
  jets ($\ctau = 10$~mm); the displaced pileup tail from strange-hadron
  decays survives the selection but is too small and uncorrelated to
  affect the observables.}
  \label{fig:pileup}
\end{figure}

\stub{In-situ negative-tail calibration (signed-$\dnought$ variants)
under each scenario.}

% =====================================================================
\section{Displacement-aware anomaly detection}
\label{sec:anomaly}

The central application of the observable family is to extend the
generality of substructure-based anomaly detection to the lifetime
frontier. The proposed protocol is deliberately minimal, so that the
comparison isolates the input basis rather than the architecture:

\begin{enumerate}
  \item \textbf{Background model.} Train a (variational) autoencoder on
  QCD jets only (signal-free training sample; in data this would be a
  control region or the bulk of an inclusive jet sample), with a
  per-jet feature vector as input.
  \item \textbf{Two input bases.} Basis $S$ (``standard''): the
  nominal column of Table~\ref{tab:observables} --- angularities,
  $\mathrm{ECF}_n$ and ratios, $\tau_N$ and ratios, $\pT^D$,
  $n_{\mathrm{trk}}$, jet mass. Basis $S{+}D$: the same vector
  augmented with the displacement-weighted column (anchor weighted
  angularities, dEEC with both pairings, low lifetime moments and
  ratio).
  \item \textbf{Evaluation.} Anomaly score = reconstruction error (AE)
  or a combination with the KL term (VAE). Compare ROC curves of the
  anomaly score against emerging-jet signal, per $\ctau$ point and per
  background flavor, for $S$ versus $S{+}D$. The figure of merit is
  the improvement in signal efficiency at fixed background anomaly
  rate (e.g.\ $10^{-3}$), as a function of $\ctau$.
  \item \textbf{Background-model stability.} Verify that augmenting
  the basis does not degrade the background model itself: the QCD
  anomaly-score distribution and its flavor composition should be
  statistically unchanged between $S$ and $S{+}D$ trainings, since for
  background the added observables are concentrated near the prompt
  pedestal $w \simeq \ln 2$ (with $b$ jets populating a mild tail ---
  itself a useful in-situ validation region).
\end{enumerate}

The expectation from Sec.~\ref{sec:results} is unambiguous: since the
supervised separation in basis $S$ is bounded by rejections of
$2$--$5$ while basis $S{+}D$ supports rejections of
$\mathcal{O}(10^2$--$10^3)$, the unsupervised ROC must improve
dramatically once displacement enters the input basis; the
quantitative question is how much of the supervised gain the
autoencoder realizes, and how stably across $\ctau$.

Figures~\ref{fig:vaescore} and~\ref{fig:vae} show the result at truth
level, for a fully-connected (V)AE trained on inclusive QCD jets only
(60/20/20 train/validation/test split; anomaly thresholds set on the
QCD test sample). The anomaly-score distributions
(Fig.~\ref{fig:vaescore}) display the mechanism directly: in basis $S$
the signal score distribution is indistinguishable from QCD, while in
$S{+}D$ the signal population migrates orders of magnitude into the
reconstruction-error tail.

\begin{figure}[tb]
  \centering
  \includegraphics[width=0.95\textwidth]{vae_scores_truth.png}
  \caption{Anomaly-score distributions (per-sample loss) for the QCD
  test sample and signal at $\ctau = 10$~mm, for the nominal basis $S$
  (left) and the augmented basis $S{+}D$ (right), AE and VAE.
  Truth-level tracking.}
  \label{fig:vaescore}
\end{figure} The blind spot is total: in basis $S$ the anomaly
detector tags emerging jets at $0.5$--$2\%$ efficiency for a $1\%$ QCD
anomaly rate --- consistent with flagging signal at the background rate
--- \emph{independent of lifetime}, and at a $10^{-3}$ anomaly rate the
efficiency is zero to sample precision. Augmenting the input basis to
$S{+}D$, with the architecture, training, and working points unchanged,
lifts the efficiency at the $1\%$ working point to $54\%$ at
$\ctau = 1$~mm, $76\%$ at 3~mm, $90\%$ at 10~mm, and $\geq 97\%$ for
$\ctau \geq 30$~mm (VAE). At these training-sample sizes the
variational architecture clearly outperforms the deterministic AE at
short lifetime ($54\%$ versus $23\%$ at 1~mm): the KL term regularizes
the background model against the heavy-tailed displaced features that
otherwise dominate the plain MSE objective. The background model is
stable under the augmentation: the $b$-jet anomaly rate at the $1\%$
inclusive working point rises only from $\sim 0.4\%$ to $\sim 1.5\%$,
the anticipated mild heavy-flavor tail that furnishes an in-situ
validation region.

Figure~\ref{fig:vae} also shows the natural ceiling for this
comparison: a \emph{supervised} classifier (gradient-boosted tree) on
the same $S{+}D$ basis, trained per $\ctau$ point and evaluated at the
identical fixed-anomaly-rate working points. Since basis and working
point are shared, the gap between the black and red curves isolates the
cost of model-independence alone. That cost is confined to short
lifetimes: at the $1\%$ working point the unsupervised VAE realizes
$54\%$ of the supervised efficiency at $\ctau = 1$~mm, $77\%$ at 3~mm,
$91\%$ at 10~mm, and $\geq 98\%$ for $\ctau \geq 30$~mm --- at long
lifetime the signal is so deep in the tails of the learned QCD density
that no supervision is needed, while near the $B$-hadron lifetime a
targeted search retains a genuine advantage over the model-agnostic
one.

\begin{figure}[tb]
  \centering
  \includegraphics[width=0.85\textwidth]{vae_efficiency_truth.png}
  \caption{Signal efficiency of the (V)AE anomaly score at a fixed QCD
  anomaly rate of $10^{-2}$ (left) and $10^{-3}$ (right), versus
  $\ctau(\pid)$. Gray dashed: standard substructure basis $S$; red
  solid: augmented basis $S{+}D$; black: supervised BDT on the same
  $S{+}D$ basis at the same working points, the per-basis ceiling that
  isolates the cost of model-independence. Identical (V)AE architecture
  and training in all cases; truth-level tracking.}
  \label{fig:vae}
\end{figure}

One design choice deserves an explicit warning, because the natural
instinct backfires. Several displaced observables are heavy-tailed
($L_1$, $\langle|\dnought|/\sigma\rangle$ span decades), and standard
autoencoder practice recommends compressing such features (e.g.\
$x \to \ln(1+x)$) before z-score normalization so the MSE objective is
not dominated by tail events. We tested exactly this and find it
\emph{degrades} the anomaly detection substantially: the VAE efficiency
at the $1\%$ working point drops from $54\%$ to $20\%$ at
$\ctau = 1$~mm and from $90\%$ to $59\%$ at 10~mm. The mechanism is
that the tails are the signal: an untransformed displaced feature
sitting many standard deviations from the QCD mean produces an
enormous reconstruction error, and compressing the tails improves
background-model fidelity at the direct expense of exactly that
anomaly contrast. For displacement-aware anomaly detection the
heavy-tailed features should enter the autoencoder untransformed.

\stub{Repeat under the standard and standard+LRT tracking scenarios;
per-flavor efficiency curves; latent-space diagnostics; cross-check
with the reco-level set-based architectures (deep-set / DETR) of the
companion framework.}

% =====================================================================
\section{The interpretability gap}
\label{sec:gap}

The ceiling is a fully supervised XGBoost classifier on the complete
$S{+}D$ observable basis, trained per $\ctau$ point (signal versus
inclusive QCD, class-weighted, early-stopped on a held-out validation
split). Against it, on identical jet splits, we compare the tree
restricted to the nominal basis $S$ and the strongest single
observables --- $\mathrm{dEEC}(\min)$ and the weighted multiplicity
$\Sigma_i w_i$. The question the comparison answers is how much of the
family's combined discrimination a single interpretable observable
retains, and how much is recoverable without lifetime information at
all.

Figure~\ref{fig:gap} shows the outcome. The single observables capture
$50$--$85\%$ of the ceiling's log-rejection across the lifetime grid,
with complementary strengths: the dEEC leads below $\sim 10$~mm, the
weighted multiplicity above (at 300~mm its measured rejection even
exceeds the test-sample-saturated ceiling bound, an artifact of the
ceiling's statistics cap). The nominal-basis tree is the supervised
counterpart of the anomaly-detection blind spot: supervised training
does extract a lifetime-\emph{independent} $\sim 50\times$ $b$-jet
rejection from shape information alone, but the curve never rises with
$\ctau$ --- displacement is where all the lifetime discrimination
lives, and two interpretable observables transmit most of it.
\todo{seed ensembles for error bands; production statistics}

\begin{figure}[tb]
  \centering
  \includegraphics[width=0.95\textwidth]{interpretability_gap_truth.png}
  \caption{Left: $b$-jet rejection at 50\% signal efficiency versus
  $\ctau(\pid)$ for the supervised XGBoost ceiling on the full $S{+}D$
  basis, the tree restricted to the nominal basis $S$, and the two
  strongest single observables. Plateaus reflect the test-sample
  statistics. Right: fraction of the ceiling's log-rejection captured
  by each proxy. Truth-level tracking.}
  \label{fig:gap}
\end{figure}

% =====================================================================
\section{Conclusions}
\label{sec:conclusions}

\stub{Summarize: substructure-based anomaly detection is blind to
displacement, and a small interpretable observable set closes the blind
spot; dEEC as the key observable at short lifetime; (V)AE ROC gain from
the augmented basis; X\% of the supervised ceiling; robustness
statements; outlook paragraph (other portals, SUEP regime,
trigger-level application, event-level anomaly detection).}

% =====================================================================
\appendix

\section{Alternate mediator portal}
\label{app:portal}
\stub{One alternate portal (e.g.\ $t$-channel bifundamental as in the
original emerging-jets model), MadGraph + Pythia HV chain; show the
observable family is portal-independent.}

\section{Pythia Hidden Valley settings}
\label{app:settings}
\todo{full card listing for reproducibility (cards/pythia/ in the code
release)}

% =====================================================================
\begin{thebibliography}{99}

\bibitem{Schwaller:2015gea}
P.~Schwaller, D.~Stolarski and A.~Weiler,
\emph{Emerging Jets},
JHEP \textbf{05} (2015) 059, [\texttt{arXiv:1502.05409}].

\bibitem{Albouy:2022cin}
G.~Albouy et al.,
\emph{Theory, phenomenology, and experimental avenues for dark showers:
a Snowmass 2021 report},
Eur.\ Phys.\ J.\ C \textbf{82} (2022) 1132, [\texttt{arXiv:2203.09503}].

\bibitem{ATLAS:EJ}
ATLAS Collaboration,
\emph{Search for emerging jets},
\todo{cite the ATLAS emerging-jets search}.

\bibitem{CMS:EJ}
CMS Collaboration,
\emph{Search for new particles decaying to a jet and an emerging jet},
JHEP \textbf{02} (2019) 179, [\texttt{arXiv:1810.10069}].

\bibitem{Farina:2018fyg}
M.~Farina, Y.~Nakai and D.~Shih,
\emph{Searching for New Physics with Deep Autoencoders},
Phys.\ Rev.\ D \textbf{101} (2020) 075021, [\texttt{arXiv:1808.08992}].

\bibitem{Heimel:2018mkt}
T.~Heimel, G.~Kasieczka, T.~Plehn and J.~M.~Thompson,
\emph{QCD or What?},
SciPost Phys.\ \textbf{6} (2019) 030, [\texttt{arXiv:1808.08979}].

\bibitem{Kasieczka:2021xcg}
G.~Kasieczka et al.,
\emph{The LHC Olympics 2020: A Community Challenge for Anomaly
Detection in High Energy Physics},
Rept.\ Prog.\ Phys.\ \textbf{84} (2021) 124201,
[\texttt{arXiv:2101.08320}].

\bibitem{ATLAS:dijetAD}
ATLAS Collaboration,
\emph{Dijet resonance search with weak supervision using
$\sqrt{s}=13$~TeV pp collisions},
Phys.\ Rev.\ Lett.\ \textbf{125} (2020) 131801,
[\texttt{arXiv:2005.02983}].

\bibitem{Chang:2013rca}
H.-M.~Chang, M.~Procura, J.~Thaler and W.~J.~Waalewijn,
\emph{Calculating Track-Based Observables for the LHC},
Phys.\ Rev.\ Lett.\ \textbf{111} (2013) 102002,
[\texttt{arXiv:1303.6637}].

\bibitem{Li:2021zcf}
Y.~Li, I.~Moult, S.~S.~van Velzen, W.~J.~Waalewijn and H.~X.~Zhu,
\emph{Extending Precision Perturbative QCD with Track Functions},
Phys.\ Rev.\ Lett.\ \textbf{128} (2022) 182001,
[\texttt{arXiv:2108.01674}].

\bibitem{Lee:2023npz}
\todo{energy correlators on tracks reference (Moult et al.)}

\bibitem{Cacciari:2008gp}
M.~Cacciari, G.~P.~Salam and G.~Soyez,
\emph{The anti-$k_t$ jet clustering algorithm},
JHEP \textbf{04} (2008) 063, [\texttt{arXiv:0802.1189}].

\bibitem{Larkoski:2014pca}
A.~J.~Larkoski, J.~Thaler and W.~J.~Waalewijn,
\emph{Gaining (Mutual) Information about Quark/Gluon Discrimination},
JHEP \textbf{11} (2014) 129, [\texttt{arXiv:1408.3122}].

\bibitem{Larkoski:2013eya}
A.~J.~Larkoski, G.~P.~Salam and J.~Thaler,
\emph{Energy Correlation Functions for Jet Substructure},
JHEP \textbf{06} (2013) 108, [\texttt{arXiv:1305.0007}].

\bibitem{Thaler:2010tr}
J.~Thaler and K.~Van Tilburg,
\emph{Identifying Boosted Objects with N-subjettiness},
JHEP \textbf{03} (2011) 015, [\texttt{arXiv:1011.2268}].

\bibitem{Bierlich:2022pfr}
C.~Bierlich et al.,
\emph{A comprehensive guide to the physics and usage of PYTHIA 8.3},
SciPost Phys.\ Codeb.\ \textbf{2022} (2022) 8,
[\texttt{arXiv:2203.11601}].

\bibitem{Carloni:2010tw}
L.~Carloni, J.~Rathsman and T.~Sj\"ostrand,
\emph{Discerning Secluded Sector gauge structures},
JHEP \textbf{04} (2011) 091, [\texttt{arXiv:1102.3795}].

\end{thebibliography}

\end{document}
